%----------------------------------------------------------------------------
\chapter{System Design}
\label{chap:system-design}
%----------------------------------------------------------------------------

\section{Assumptions}
\label{sec:system-design-assumptions}

The localization and mapping system developed in this thesis is built on a
small set of modeling assumptions. Their purpose is to keep the estimation
problem tractable for a simultaneous localization and mapping (SLAM)
formulation based on the extended Kalman filter (EKF), while remaining
physically justified for an autonomous surface vehicle (ASV) operating on a
sheltered inland course such as the RoboBoat competition venue. Each
assumption is stated below together with its justification and the
consequences of its violation. The two central assumptions, planar motion
and a static environment, are discussed in
Sections~\ref{sec:assumption-planar} and~\ref{sec:assumption-static};
further assumptions concerning the sensors and the estimator are collected
in Section~\ref{sec:assumption-further}.

%----------------------------------------------------------------------------
\subsection{Planar Motion}
\label{sec:assumption-planar}
%----------------------------------------------------------------------------

The vehicle is assumed to move in the horizontal plane of a local,
Earth-fixed navigation frame $\mathcal{W}$, whose axes point east, north and
up (ENU) in accordance with the ROS coordinate conventions~\cite{rep103}.
The pose of the vehicle body frame $\mathcal{B}$ at discrete time step $k$
is described by
\begin{equation}
  \mathbf{x}_k =
  \begin{bmatrix} x_k & y_k & \psi_k \end{bmatrix}^{\top},
  \qquad \psi_k \in (-\pi, \pi],
  \label{eq:planar-pose}
\end{equation}
where $(x_k, y_k)$ is the position of the origin of $\mathcal{B}$ in the
east--north plane and $\psi_k$ is the yaw angle, measured counter-clockwise
from the east axis. Equivalently, the pose is an element of the special
Euclidean group $\mathrm{SE}(2)$ with the homogeneous representation
\begin{equation}
  \mathbf{T}_k =
  \begin{bmatrix}
    \cos\psi_k & -\sin\psi_k & x_k \\
    \sin\psi_k & \cos\psi_k  & y_k \\
    0          & 0           & 1
  \end{bmatrix} \in \mathrm{SE}(2).
  \label{eq:planar-se2}
\end{equation}
% TODO: If the process model carries velocities or IMU biases as states,
% extend \mathbf{x}_k here and refer to the process model section.

The restriction to three degrees of freedom is justified by the following
observations.

\begin{itemize}
  \item An ASV is constrained to the water surface. Its heave, roll and
  pitch are bounded, oscillatory responses to waves and maneuvers; they
  carry no information relevant for navigation, and all competition tasks
  are defined in the horizontal plane.

  \item The remaining degrees of freedom are not ignored but treated as
  known inputs. The roll angle $\phi_k$ and pitch angle $\theta_k$, using
  the Z-Y-X (yaw--pitch--roll) Euler angle convention, are provided by the
  attitude estimate of the flight controller's inertial measurement unit
  (IMU). They are used to rotate every LiDAR and camera measurement into a
  gravity-aligned frame before it is projected onto the horizontal plane:
  \begin{equation}
    \mathbf{p}^{\mathcal{H}} =
    \mathbf{R}_y(\theta_k)\,\mathbf{R}_x(\phi_k)
    \left( \mathbf{R}_{\mathcal{B}\mathcal{S}}\,\mathbf{p}^{\mathcal{S}}
         + \mathbf{t}_{\mathcal{B}\mathcal{S}} \right),
    \label{eq:tilt-compensation}
  \end{equation}
  where $\mathbf{p}^{\mathcal{S}}$ is a point in the sensor frame
  $\mathcal{S}$,
  $(\mathbf{R}_{\mathcal{B}\mathcal{S}}, \mathbf{t}_{\mathcal{B}\mathcal{S}})$
  is the sensor-to-body transform, and $\mathcal{H}$ denotes the level frame
  that shares the heading of $\mathcal{B}$. Only the first two components
  of $\mathbf{p}^{\mathcal{H}}$ are used.

  \item The operating area extends over at most a few hundred meters. GNSS
  positions are converted from geodetic coordinates to the local ENU frame,
  whose origin is fixed at a surveyed reference point, and the up component
  is discarded. For a point at a distance $s$ along the Earth's surface from
  the origin, the horizontal ENU coordinates deviate from the true distance
  by approximately $s^3 / (6 R_\mathrm{E}^2)$, where
  $R_\mathrm{E} \approx 6371\,\mathrm{km}$. At $s = 1\,\mathrm{km}$ this
  amounts to about $4\,\mu\mathrm{m}$, which is negligible.

  \item The heading measured by the dual-antenna GNSS receiver is referenced
  to geodetic north, i.e.\ to the same frame as the position measurements,
  so no magnetic declination has to be modeled. The azimuth $\alpha$ of the
  vehicle's longitudinal axis, measured clockwise from north after
  correcting for the orientation of the antenna baseline
  (cf.~\ref{asm:extrinsics}), is converted to the yaw angle of
  \eqref{eq:planar-pose} by $\psi = \pi/2 - \alpha$, followed by wrapping
  to $(-\pi, \pi]$.
\end{itemize}

As a consequence, attitude errors only enter the planar problem through the
projection~\eqref{eq:tilt-compensation}. Without compensation, a tilt
$\theta$ in the direction of a target at sensor height would shorten its
measured horizontal range by $r(1-\cos\theta) \approx r\theta^2/2$, e.g.\
by $7.6\,\mathrm{cm}$ at $r = 20\,\mathrm{m}$ and $\theta = 5^\circ$, and
targets above or below the sensor would additionally be displaced by
$\Delta h \tan\theta$. With the IMU attitude applied, only the attitude
estimation error remains, which is small compared with the other error
sources.
% TODO: Confirm typical roll/pitch amplitudes from the simulation logs.

%----------------------------------------------------------------------------
\subsection{Static Landmarks}
\label{sec:assumption-static}
%----------------------------------------------------------------------------

The map is represented by $N$ point landmarks
\begin{equation}
  \mathbf{m}_j =
  \begin{bmatrix} m_{x,j} & m_{y,j} \end{bmatrix}^{\top} \in \mathbb{R}^2,
  \qquad j = 1, \dots, N,
  \label{eq:landmark}
\end{equation}
each corresponding to a moored course marker, primarily the navigation
buoys of the competition course. The landmarks are assumed to be static in
$\mathcal{W}$,
\begin{equation}
  \mathbf{m}_{j,k+1} = \mathbf{m}_{j,k}, \qquad j = 1, \dots, N,
  \label{eq:static-landmarks}
\end{equation}
so that the full state estimated by the filter is
\begin{equation}
  \boldsymbol{\xi}_k =
  \begin{bmatrix}
    \mathbf{x}_k^{\top} & \mathbf{m}_1^{\top} & \cdots & \mathbf{m}_N^{\top}
  \end{bmatrix}^{\top} \in \mathbb{R}^{3+2N}.
  \label{eq:slam-state}
\end{equation}

This assumption reflects how the course is built: buoys are anchored, and
their layout does not change during a run. Three deviations from the ideal
are expected and handled as follows.

\begin{itemize}
  \item \textbf{Mooring motion.} A moored buoy moves within a small circle
  around its anchor under the influence of wind, current and waves. This
  motion is small compared with the other error sources and is absorbed by
  the landmark measurement noise covariance rather than modeled as a
  separate process.
  % TODO: Quantify the expected buoy excursion if data is available.

  \item \textbf{Finite extent.} A buoy of radius $r_\mathrm{b}$ is not a
  point, and the LiDAR only observes the side facing the sensor. For a
  circular cross-section sampled uniformly in angle from a distant sensor,
  the centroid of the returns lies approximately $\pi r_\mathrm{b} / 4$
  closer to the sensor than the buoy's axis, e.g.\ about $0.12\,\mathrm{m}$
  for $r_\mathrm{b} = 0.15\,\mathrm{m}$. Since the camera classifies each
  buoy, this bias can be compensated with a class-dependent offset along
  the line of sight.

  \item \textbf{Dynamic objects.} Other vessels, safety boats, wildlife and
  the vehicle's own wake violate the static assumption. They are kept out
  of the map: only detections classified as a landmark class are admitted
  as landmark candidates, while all other detections are used by the
  obstacle avoidance layer only, outside the SLAM state.
\end{itemize}

The static assumption is essential to the behavior of EKF-SLAM. For a
static map, the determinant of any submatrix of the map covariance matrix
decreases monotonically as successive observations are
made~\cite{dissanayake2001}. A landmark that actually moves would therefore
cause an overconfident estimate and, through the correlations between
landmark and vehicle states, would also corrupt the pose estimate.
Measurements that are inconsistent with the static assumption must
consequently be rejected before the update step.

%----------------------------------------------------------------------------
\subsection{Further Assumptions}
\label{sec:assumption-further}
%----------------------------------------------------------------------------

% Requires \usepackage{enumitem} in the preamble.
\begin{enumerate}[label=\textbf{A\arabic*}, ref=A\arabic*, leftmargin=*]
  \item\label{asm:extrinsics}
  \textbf{Rigid mounting and known extrinsics.} All sensors are rigidly
  attached to the vehicle body, and the transforms between each sensor
  frame and $\mathcal{B}$ are constant and known. This includes the lever
  arms of the two GNSS antennas and the upside-down mounting of the LiDAR,
  which corresponds to a rotation of $180^\circ$ about the sensor's
  longitudinal axis. In simulation the transforms are known exactly from
  the vehicle model; on the physical vehicle they have to be obtained from
  the CAD model and refined by calibration.

  \item\label{asm:time}
  \textbf{Common time base.} All measurements are timestamped in a common
  clock, and transport latencies are either negligible or compensated, so
  that measurements can be processed in timestamp order. In simulation all
  sensors are stamped with the simulation clock; on the physical vehicle,
  hardware time synchronization (e.g.\ the Precision Time Protocol
  supported by the LiDAR) is required.

  \item\label{asm:noise}
  \textbf{Gaussian white noise.} Measurement and process noise are
  zero-mean, Gaussian, white and mutually independent between sensors, with
  known covariances derived from the sensor specifications
  (Section~\ref{sec:system-design-sensors}). Together
  with~\eqref{eq:static-landmarks}, this corresponds to the Markov
  assumption underlying recursive Bayesian estimation~\cite{thrun2005}.
  % TODO: Create sec:system-design-sensors or change this reference.

  \item\label{asm:gnss}
  \textbf{GNSS availability.} GNSS position and dual-antenna heading are
  available in the open-sky environment of the course, with an accuracy
  that depends on the correction status (RTK fixed, RTK float or
  standalone). This status is reflected in the measurement covariance
  rather than assumed constant.
  % TODO: State here whether GNSS outages are simulated in the evaluation.

  \item\label{asm:heading}
  \textbf{Small heading uncertainty.} The heading error of the vehicle
  remains small throughout operation, since heading is directly observed by
  the dual-antenna GNSS receiver. This is relevant because the
  inconsistency of EKF-SLAM has been shown to stem primarily from
  linearization errors caused by heading uncertainty~\cite{bailey2006}.

  \item\label{asm:map-size}
  \textbf{Moderate map size.} The number of landmarks is on the order of
  tens. Since the computational cost of an EKF update grows quadratically
  with the state dimension, a state of, for example,
  $3 + 2 \cdot 50 = 103$ elements remains well within real-time limits, and
  no submapping or sparsification is required.
  % TODO: Confirm the expected landmark count from the course layout.

  \item\label{asm:association}
  \textbf{Unknown correspondences.} The correspondence between measurements
  and landmarks is \emph{not} assumed to be known; it is estimated online
  from the landmark class provided by the camera and a gating test on the
  Mahalanobis distance of the innovation.
  % TODO: Adjust if known data association is assumed in simulation.
\end{enumerate}

\section{Sensor Configuration}
\label{sec:system-design-sensor-configuration}
% TODO: LiDAR/camera resolution analysis

\section{Coordinate Frames and Transformations}
\label{sec:system-design-coordinate-frames}
% TODO: TF tree, WGS84 to ENU

\section{Vehicle Motion Model}
\label{sec:system-design-motion-model}
% TODO: Derive the vehicle motion model

\section{Observation Models}
\label{sec:system-design-observation-models}
% TODO: Range-bearing landmark model, GNSS model

\section{EKF-SLAM Formulation}
\label{sec:system-design-ekf-slam-formulation}
% TODO: State vector, prediction, correction, state augmentation

\section{Noise Models}
\label{sec:system-design-noise-models}
% TODO: Datasheet-based Q and R

\section{Data Association and Map Management Design}
\label{sec:system-design-data-association}
% TODO: Data association and map management design
